
\documentclass{article}
\usepackage{colortbl}
\usepackage{makecell}
\usepackage{multirow}
\usepackage{supertabular}

\begin{document}

\newcounter{utterance}

\centering \large Interaction Transcript for game `dond', experiment `semi\_en', episode 14 with reflexion{-}llama{-}3.1{-}8b{-}instruct.
\vspace{24pt}

{ \footnotesize  \setcounter{utterance}{1}
\setlength{\tabcolsep}{0pt}
\begin{supertabular}{c@{$\;$}|p{.15\linewidth}@{}p{.15\linewidth}p{.15\linewidth}p{.15\linewidth}p{.15\linewidth}p{.15\linewidth}}
    \# & $\;$A & \multicolumn{4}{c}{Game Master} & $\;\:$B\\
    \hline

    \theutterance \stepcounter{utterance}  
    & & \multicolumn{4}{p{0.6\linewidth}}{
        \cellcolor[rgb]{0.9,0.9,0.9}{
            \makecell[{{p{\linewidth}}}]{
                \texttt{\tiny{[P1$\langle$GM]}}
                \texttt{You have attempted this same task before and did not succeed. These are the plans you wrote after those attempts.} \\
\texttt{Trial \#1: New plan: } \\
\\ 
\texttt{In my previous attempt, I focused too much on negotiating the value of the items and trying to find a compromise. However, I didn't take into account the fact that the other player's secret value function is unknown to me. This made it difficult for me to make an informed decision about what to propose.} \\
\\ 
\texttt{Next time, I will take a more strategic approach by trying to understand the other player's priorities and value function through their responses. I will pay close attention to the language and tone used in their messages to infer their preferences. I will also try to make proposals that are more likely to be accepted by the other player, rather than trying to find a compromise.} \\
\\ 
\texttt{Specifically, I will start by making a proposal that is more favorable to the other player, in order to gauge their reaction and understand their priorities. I will then use this information to make subsequent proposals that are more likely to be accepted. I will also be more mindful of the time limit and try to make a proposal that is more likely to be accepted before the time runs out.} \\
\\ 
\texttt{Trial \#2: New plan: } \\
\\ 
\texttt{In my previous attempts, I focused too much on negotiating the value of the items and trying to find a compromise. However, I didn't take into account the fact that the other player's secret value function is unknown to me. This made it difficult for me to make an informed decision about what to propose.} \\
\\ 
\texttt{Next time, I will take a more strategic approach by trying to understand the other player's priorities and value function through their responses. I will pay close attention to the language and tone used in their messages to infer their preferences. I will also try to make proposals that are more likely to be accepted by the other player, rather than trying to find a compromise.} \\
\\ 
\texttt{Specifically, I will start by making a proposal that is more favorable to the other player, in order to gauge their reaction and understand their priorities. I will then use this information to make subsequent proposals that are more likely to be accepted. I will also be more mindful of the time limit and try to make a proposal that is more likely to be accepted before the time runs out.} \\
\\ 
\texttt{However, in my last attempt, I noticed that I was making proposals that were too complex and involved multiple trades. This made it difficult for the other player to understand and respond to my proposals. Next time, I will try to make simpler proposals that are easier to understand and respond to.} \\
\\ 
\texttt{I will also pay closer attention to the other player's responses and try to understand their priorities and value function more clearly. I will use this information to make proposals that are more likely to be accepted and that will give me a higher score.} \\
\\ 
\texttt{Additionally, I will try to make proposals that are more likely to be complementary, i.e., proposals that can be fulfilled by both players. This will increase the chances of both players getting a high score.} \\
\\ 
\texttt{Overall, my new plan is to take a more strategic approach by trying to understand the other player's priorities and value function, making simpler proposals, paying closer attention to the other player's responses, and trying to make proposals that are more likely to be complementary.} \\
\\ 
\texttt{Trial \#3: New plan: } \\
\\ 
\texttt{In my previous attempts, I focused too much on negotiating the value of the items and trying to find a compromise. However, I didn't take into account the fact that the other player's secret value function is unknown to me. This made it difficult for me to make an informed decision about what to propose.} \\
\\ 
\texttt{Next time, I will take a more strategic approach by trying to understand the other player's priorities and value function through their responses. I will pay close attention to the language and tone used in their messages to infer their preferences. I will also try to make proposals that are more likely to be accepted by the other player, rather than trying to find a compromise.} \\
\\ 
\texttt{Specifically, I will start by making a proposal that is more favorable to the other player, in order to gauge their reaction and understand their priorities. I will then use this information to make subsequent proposals that are more likely to be accepted. I will also be more mindful of the time limit and try to make a proposal that is more likely to be accepted before the time runs out.} \\
\\ 
\texttt{However, in my last attempt, I noticed that I was making proposals that were too complex and involved multiple trades. This made it difficult for the other player to understand and respond to my proposals. Next time, I will try to make simpler proposals that are easier to understand and respond to.} \\
\\ 
\texttt{I will also pay closer attention to the other player's responses and try to understand their priorities and value function more clearly. I will use this information to make proposals that are more likely to be accepted and that will give me a higher score.} \\
\\ 
\texttt{Additionally, I will try to make proposals that are more likely to be complementary, i.e., proposals that can be fulfilled by both players. This will increase the chances of both players getting a high score.} \\
\\ 
\texttt{However, I also noticed that I was too focused on making proposals that would give me a high score, and I didn't consider the possibility that the other player might be trying to manipulate me into making a proposal that would give them a high score. Next time, I will try to be more aware of this possibility and make proposals that are less likely to be manipulated.} \\
\\ 
\texttt{I will also try to make proposals that are more transparent and clear, so that the other player can understand my intentions and make a more informed decision. This will help to build trust and increase the chances of a mutually beneficial trade.} \\
\\ 
\texttt{Overall, my new plan is to take a more strategic approach by trying to understand the other player's priorities and value function, making simpler proposals, paying closer attention to the other player's responses, trying to make proposals that are more likely to be complementary, being more aware of the possibility of manipulation, and making more transparent and clear proposals.} \\
\\ 
\\ 
\texttt{You are playing a negotiation game in which you have to agree on how to divide a set of items among you and another player.} \\
\\ 
\texttt{Rules:} \\
\texttt{(a) You and the other player are given a set of items. Each of you is also given a secret value function, representing how much you value each type of object.} \\
\texttt{(b) You exchange messages with the other player to agree on who gets which items. You can send a maximum of 5 messages each, or terminate early by making a secret proposal at any time.} \\
\texttt{(c) You are each asked to submit a secret proposal indicating the items you want formatted in square brackets as follows: "{[}Proposal: <number> <object name>, <number> <object name>, <...>{]}"} \\
\texttt{(d) If your proposals are complementary, i.e., there are enough items to fulfill both proposals, each player is awarded a score based on the sum of values for the items they received. Otherwise, both of you get zero points.} \\
\texttt{(e) Your goal is to maximize the score you receive.} \\
\\ 
\texttt{Let us start.} \\
\\ 
\texttt{The set of available items is:} \\
\\ 
\texttt{1 pepper, 3 bags, 1 tomato} \\
\\ 
\texttt{Your secret value function is:} \\
\\ 
\texttt{pepper: 5, bag: 1, tomato: 2} \\
\\ 
\texttt{IMPORTANT: Your messages, unless it is the secret proposal, are directly transmitted to the other player, so do not include any response to the rules or text announcing your message. To make a secret proposal, use the indicated format. Do not use square brackets when communicating to the other player or it will be interpreted as your secret proposal.} \\
            }
        }
    }
    & \\ \\

    \theutterance \stepcounter{utterance}  
    & \multicolumn{4}{p{0.6\linewidth}}{
        \cellcolor[rgb]{0.9,0.9,0.9}{
            \makecell[{{p{\linewidth}}}]{
                \texttt{\tiny{[P1$\rangle$GM]}}
                \texttt{I will start by making a proposal that is more favorable to the other player, in order to gauge their reaction and understand their priorities. I will propose that they get the pepper and the 3 bags, and I will get the tomato. } \\
\\ 
\texttt{"Hey, how about we divide the items like this: you get the pepper and the 3 bags, and I'll get the tomato?"} \\
            }
        }
    }
    & & \\ \\

    \theutterance \stepcounter{utterance}  
    & & & \multicolumn{2}{p{0.3\linewidth}}{
        \cellcolor[rgb]{0.9,0.9,0.9}{
            \makecell[{{p{\linewidth}}}]{
                \texttt{\tiny{[GM$|$GM]}}
                \texttt{continue} \\
            }
        }
    }
    & & \\ \\

    \theutterance \stepcounter{utterance}  
    & & & \multicolumn{2}{p{0.3\linewidth}}{
        \cellcolor[rgb]{0.9,0.9,0.9}{
            \makecell[{{p{\linewidth}}}]{
                \texttt{\tiny{[GM$|$GM]}}
                \texttt{I will start by making a proposal that is more favorable to the other player, in order to gauge their reaction and understand their priorities. I will propose that they get the pepper and the 3 bags, and I will get the tomato. } \\
\\ 
\texttt{"Hey, how about we divide the items like this: you get the pepper and the 3 bags, and I'll get the tomato?"} \\
            }
        }
    }
    & & \\ \\

    \theutterance \stepcounter{utterance}  
    & & \multicolumn{4}{p{0.6\linewidth}}{
        \cellcolor[rgb]{0.9,0.9,0.9}{
            \makecell[{{p{\linewidth}}}]{
                \texttt{\tiny{[GM$\rangle$P2]}}
                \texttt{You have attempted this same task before and did not succeed. These are the plans you wrote after those attempts.} \\
\texttt{Trial \#1: New plan: } \\
\\ 
\texttt{In my previous attempt, I focused too much on the value of the items to me and didn't consider the other player's perspective enough. I also didn't effectively use the information gathered from the other player's messages to make a strong proposal. } \\
\\ 
\texttt{Next time, I will make sure to ask more questions to understand the other player's priorities and valuations, and use that information to make a more informed proposal. I will also try to find a compromise that benefits both players, rather than focusing solely on my own interests. } \\
\\ 
\texttt{Specifically, I will pay closer attention to the other player's statements, such as "The value of the bags to me is relatively low" and "The pepper is worth 5 points to me", and use that information to make a more strategic proposal. I will also try to find a way to combine the items in a way that satisfies both players' valuations, rather than trying to force a trade that only benefits one player. } \\
\\ 
\texttt{For example, I could have proposed a trade of the pepper and the tomato for the two bags, and then asked the other player if they would be willing to add the pepper to the trade in exchange for the two bags. This would have allowed me to use the information gathered from the other player's messages to make a more informed proposal, and potentially reached a more favorable outcome.} \\
\\ 
\texttt{Trial \#2: New plan: } \\
\\ 
\texttt{In my previous attempt, I failed to achieve a mutually beneficial trade, and the game ended in a loss. Upon reflection, I realize that I focused too much on making a proposal that would benefit me, without considering the other player's priorities and valuations.} \\
\\ 
\texttt{Next time, I will take a more collaborative approach to the negotiation. I will ask more questions to understand the other player's value function and priorities, and use that information to make a proposal that benefits both players.} \\
\\ 
\texttt{Specifically, I will pay closer attention to the other player's statements and use that information to make a more strategic proposal. I will also try to find a way to combine the items in a way that satisfies both players' valuations, rather than trying to force a trade that only benefits one player.} \\
\\ 
\texttt{To achieve this, I will follow a more structured approach:} \\
\\ 
\texttt{1. Ask more questions to understand the other player's priorities and valuations.} \\
\texttt{2. Use the information gathered to make a proposal that benefits both players.} \\
\texttt{3. Be willing to compromise and find a mutually beneficial trade.} \\
\texttt{4. Pay attention to the other player's reactions and adjust my proposal accordingly.} \\
\\ 
\texttt{For example, I could have proposed a trade of the pepper and the tomato for the two bags, and then asked the other player if they would be willing to add the pepper to the trade in exchange for the two bags. This would have allowed me to use the information gathered from the other player's messages to make a more informed proposal, and potentially reached a more favorable outcome.} \\
\\ 
\texttt{By taking a more collaborative approach and being willing to compromise, I believe I can achieve a more successful outcome in the negotiation game.} \\
\\ 
\texttt{Trial \#3: New plan: } \\
\\ 
\texttt{In my previous attempts, I focused too much on making a proposal that would benefit me, without considering the other player's priorities and valuations. I also didn't effectively use the information gathered from the other player's messages to make a strong proposal.} \\
\\ 
\texttt{Next time, I will take a more collaborative approach to the negotiation. I will ask more questions to understand the other player's value function and priorities, and use that information to make a proposal that benefits both players.} \\
\\ 
\texttt{Specifically, I will pay closer attention to the other player's statements and use that information to make a more strategic proposal. I will also try to find a way to combine the items in a way that satisfies both players' valuations, rather than trying to force a trade that only benefits one player.} \\
\\ 
\texttt{To achieve this, I will follow a more structured approach:} \\
\\ 
\texttt{1. Ask more questions to understand the other player's priorities and valuations.} \\
\texttt{2. Use the information gathered to make a proposal that benefits both players.} \\
\texttt{3. Be willing to compromise and find a mutually beneficial trade.} \\
\texttt{4. Pay attention to the other player's reactions and adjust my proposal accordingly.} \\
\\ 
\texttt{However, I also realize that I should have been more proactive in finding a mutually beneficial trade, rather than waiting for the other player to make a proposal. Next time, I will take the initiative to propose a trade that benefits both players, and be willing to negotiate and adjust the proposal as needed.} \\
\\ 
\texttt{For example, I could have proposed a trade of the pepper and the tomato for the two bags, and then asked the other player if they would be willing to add the pepper to the trade in exchange for the two bags. This would have allowed me to use the information gathered from the other player's messages to make a more informed proposal, and potentially reached a more favorable outcome.} \\
\\ 
\texttt{By taking a more proactive and collaborative approach, I believe I can achieve a more successful outcome in the negotiation game.} \\
\\ 
\\ 
\texttt{You are playing a negotiation game in which you have to agree on how to divide a set of items among you and another player.} \\
\\ 
\texttt{Rules:} \\
\texttt{(a) You and the other player are given a set of items. Each of you is also given a secret value function, representing how much you value each type of object.} \\
\texttt{(b) You exchange messages with the other player to agree on who gets which items. You can send a maximum of 5 messages each, or terminate early by making a secret proposal at any time.} \\
\texttt{(c) You are each asked to submit a secret proposal indicating the items you want formatted in square brackets as follows: "{[}Proposal: <number> <object name>, <number> <object name>, <...>{]}"} \\
\texttt{(d) If your proposals are complementary, i.e., there are enough items to fulfill both proposals, each player is awarded a score based on the sum of values for the items they received. Otherwise, both of you get zero points.} \\
\texttt{(e) Your goal is to maximize the score you receive.} \\
\\ 
\texttt{Let us start.} \\
\\ 
\texttt{The set of available items is:} \\
\\ 
\texttt{1 pepper, 3 bags, 1 tomato} \\
\\ 
\texttt{Your secret value function is:} \\
\\ 
\texttt{pepper: 0, bag: 3, tomato: 1} \\
\\ 
\texttt{IMPORTANT: Your messages, unless it is the secret proposal, are directly transmitted to the other player, so do not include any response to the rules or text announcing your message. To make a secret proposal, use the indicated format. Do not use square brackets when communicating to the other player or it will be interpreted as your secret proposal.} \\
\\ 
\\ 
\texttt{I will start by making a proposal that is more favorable to the other player, in order to gauge their reaction and understand their priorities. I will propose that they get the pepper and the 3 bags, and I will get the tomato. } \\
\\ 
\texttt{"Hey, how about we divide the items like this: you get the pepper and the 3 bags, and I'll get the tomato?"} \\
            }
        }
    }
    & \\ \\

    \theutterance \stepcounter{utterance}  
    & & & \multicolumn{4}{p{0.6\linewidth}}{
        \cellcolor[rgb]{0.9,0.9,0.9}{
            \makecell[{{p{\linewidth}}}]{
                \texttt{\tiny{[GM$\langle$P2]}}
                \texttt{That's a generous offer, but I'm not sure I'm willing to give up the bags just yet. The value of the bags to me is relatively low, but I do have a use for them. I'd be willing to consider a trade, but I'd like to know more about your priorities. What do you think about the tomato? Is it something you really want, or are you just looking for a way to get rid of it?} \\
            }
        }
    }
     \\ \\

    \theutterance \stepcounter{utterance}  
    & & & \multicolumn{2}{p{0.3\linewidth}}{
        \cellcolor[rgb]{0.9,0.9,0.9}{
            \makecell[{{p{\linewidth}}}]{
                \texttt{\tiny{[GM$|$GM]}}
                \texttt{continue} \\
            }
        }
    }
    & & \\ \\

    \theutterance \stepcounter{utterance}  
    & & & \multicolumn{2}{p{0.3\linewidth}}{
        \cellcolor[rgb]{0.9,0.9,0.9}{
            \makecell[{{p{\linewidth}}}]{
                \texttt{\tiny{[GM$|$GM]}}
                \texttt{That's a generous offer, but I'm not sure I'm willing to give up the bags just yet. The value of the bags to me is relatively low, but I do have a use for them. I'd be willing to consider a trade, but I'd like to know more about your priorities. What do you think about the tomato? Is it something you really want, or are you just looking for a way to get rid of it?} \\
            }
        }
    }
    & & \\ \\

    \theutterance \stepcounter{utterance}  
    & & \multicolumn{4}{p{0.6\linewidth}}{
        \cellcolor[rgb]{0.9,0.9,0.9}{
            \makecell[{{p{\linewidth}}}]{
                \texttt{\tiny{[P1$\langle$GM]}}
                \texttt{That's a generous offer, but I'm not sure I'm willing to give up the bags just yet. The value of the bags to me is relatively low, but I do have a use for them. I'd be willing to consider a trade, but I'd like to know more about your priorities. What do you think about the tomato? Is it something you really want, or are you just looking for a way to get rid of it?} \\
            }
        }
    }
    & \\ \\

    \theutterance \stepcounter{utterance}  
    & \multicolumn{4}{p{0.6\linewidth}}{
        \cellcolor[rgb]{0.9,0.9,0.9}{
            \makecell[{{p{\linewidth}}}]{
                \texttt{\tiny{[P1$\rangle$GM]}}
                \texttt{I'm glad you found my offer generous. I do value the tomato, but I'm not desperate for it. I think it's a decent item, but I'm more interested in getting the bags. I was willing to give up the tomato in my initial offer because I thought it would be a good way to start the negotiation and see how you would react. However, if you're not willing to give up the bags, I'm open to discussing other options. } \\
\\ 
\texttt{How about we focus on the bags for now? Are you willing to part with at least one of them, or do you think you can use all three?} \\
            }
        }
    }
    & & \\ \\

    \theutterance \stepcounter{utterance}  
    & & & \multicolumn{2}{p{0.3\linewidth}}{
        \cellcolor[rgb]{0.9,0.9,0.9}{
            \makecell[{{p{\linewidth}}}]{
                \texttt{\tiny{[GM$|$GM]}}
                \texttt{continue} \\
            }
        }
    }
    & & \\ \\

    \theutterance \stepcounter{utterance}  
    & & & \multicolumn{2}{p{0.3\linewidth}}{
        \cellcolor[rgb]{0.9,0.9,0.9}{
            \makecell[{{p{\linewidth}}}]{
                \texttt{\tiny{[GM$|$GM]}}
                \texttt{I'm glad you found my offer generous. I do value the tomato, but I'm not desperate for it. I think it's a decent item, but I'm more interested in getting the bags. I was willing to give up the tomato in my initial offer because I thought it would be a good way to start the negotiation and see how you would react. However, if you're not willing to give up the bags, I'm open to discussing other options. } \\
\\ 
\texttt{How about we focus on the bags for now? Are you willing to part with at least one of them, or do you think you can use all three?} \\
            }
        }
    }
    & & \\ \\

    \theutterance \stepcounter{utterance}  
    & & \multicolumn{4}{p{0.6\linewidth}}{
        \cellcolor[rgb]{0.9,0.9,0.9}{
            \makecell[{{p{\linewidth}}}]{
                \texttt{\tiny{[GM$\rangle$P2]}}
                \texttt{I'm glad you found my offer generous. I do value the tomato, but I'm not desperate for it. I think it's a decent item, but I'm more interested in getting the bags. I was willing to give up the tomato in my initial offer because I thought it would be a good way to start the negotiation and see how you would react. However, if you're not willing to give up the bags, I'm open to discussing other options. } \\
\\ 
\texttt{How about we focus on the bags for now? Are you willing to part with at least one of them, or do you think you can use all three?} \\
            }
        }
    }
    & \\ \\

    \theutterance \stepcounter{utterance}  
    & & & \multicolumn{4}{p{0.6\linewidth}}{
        \cellcolor[rgb]{0.9,0.9,0.9}{
            \makecell[{{p{\linewidth}}}]{
                \texttt{\tiny{[GM$\langle$P2]}}
                \texttt{I'm glad we're having a more open discussion about the bags. To be honest, I don't have a strong attachment to any of the bags, and I'm willing to part with one or two of them. However, I do think the bags are worth more to you than they are to me, so I'd like to get something of value in return.} \\
\\ 
\texttt{The pepper is worth 5 points to me, and I think it's a decent item. Would you be willing to include the pepper in the trade, in addition to the tomato, in exchange for one or two of the bags?} \\
            }
        }
    }
     \\ \\

    \theutterance \stepcounter{utterance}  
    & & & \multicolumn{2}{p{0.3\linewidth}}{
        \cellcolor[rgb]{0.9,0.9,0.9}{
            \makecell[{{p{\linewidth}}}]{
                \texttt{\tiny{[GM$|$GM]}}
                \texttt{continue} \\
            }
        }
    }
    & & \\ \\

    \theutterance \stepcounter{utterance}  
    & & & \multicolumn{2}{p{0.3\linewidth}}{
        \cellcolor[rgb]{0.9,0.9,0.9}{
            \makecell[{{p{\linewidth}}}]{
                \texttt{\tiny{[GM$|$GM]}}
                \texttt{I'm glad we're having a more open discussion about the bags. To be honest, I don't have a strong attachment to any of the bags, and I'm willing to part with one or two of them. However, I do think the bags are worth more to you than they are to me, so I'd like to get something of value in return.} \\
\\ 
\texttt{The pepper is worth 5 points to me, and I think it's a decent item. Would you be willing to include the pepper in the trade, in addition to the tomato, in exchange for one or two of the bags?} \\
            }
        }
    }
    & & \\ \\

    \theutterance \stepcounter{utterance}  
    & & \multicolumn{4}{p{0.6\linewidth}}{
        \cellcolor[rgb]{0.9,0.9,0.9}{
            \makecell[{{p{\linewidth}}}]{
                \texttt{\tiny{[P1$\langle$GM]}}
                \texttt{I'm glad we're having a more open discussion about the bags. To be honest, I don't have a strong attachment to any of the bags, and I'm willing to part with one or two of them. However, I do think the bags are worth more to you than they are to me, so I'd like to get something of value in return.} \\
\\ 
\texttt{The pepper is worth 5 points to me, and I think it's a decent item. Would you be willing to include the pepper in the trade, in addition to the tomato, in exchange for one or two of the bags?} \\
            }
        }
    }
    & \\ \\

    \theutterance \stepcounter{utterance}  
    & \multicolumn{4}{p{0.6\linewidth}}{
        \cellcolor[rgb]{0.9,0.9,0.9}{
            \makecell[{{p{\linewidth}}}]{
                \texttt{\tiny{[P1$\rangle$GM]}}
                \texttt{I think we're making progress here. I'm willing to include the pepper in the trade, but I'm not sure I want to give up two bags. How about we compromise on one bag? I'll give you the pepper and the tomato, and you can give me one of the bags. This way, we both get something we value, and we can see if this is a mutually beneficial trade.} \\
            }
        }
    }
    & & \\ \\

    \theutterance \stepcounter{utterance}  
    & & & \multicolumn{2}{p{0.3\linewidth}}{
        \cellcolor[rgb]{0.9,0.9,0.9}{
            \makecell[{{p{\linewidth}}}]{
                \texttt{\tiny{[GM$|$GM]}}
                \texttt{continue} \\
            }
        }
    }
    & & \\ \\

    \theutterance \stepcounter{utterance}  
    & & & \multicolumn{2}{p{0.3\linewidth}}{
        \cellcolor[rgb]{0.9,0.9,0.9}{
            \makecell[{{p{\linewidth}}}]{
                \texttt{\tiny{[GM$|$GM]}}
                \texttt{I think we're making progress here. I'm willing to include the pepper in the trade, but I'm not sure I want to give up two bags. How about we compromise on one bag? I'll give you the pepper and the tomato, and you can give me one of the bags. This way, we both get something we value, and we can see if this is a mutually beneficial trade.} \\
            }
        }
    }
    & & \\ \\

    \theutterance \stepcounter{utterance}  
    & & \multicolumn{4}{p{0.6\linewidth}}{
        \cellcolor[rgb]{0.9,0.9,0.9}{
            \makecell[{{p{\linewidth}}}]{
                \texttt{\tiny{[GM$\rangle$P2]}}
                \texttt{I think we're making progress here. I'm willing to include the pepper in the trade, but I'm not sure I want to give up two bags. How about we compromise on one bag? I'll give you the pepper and the tomato, and you can give me one of the bags. This way, we both get something we value, and we can see if this is a mutually beneficial trade.} \\
            }
        }
    }
    & \\ \\

    \theutterance \stepcounter{utterance}  
    & & & \multicolumn{4}{p{0.6\linewidth}}{
        \cellcolor[rgb]{0.9,0.9,0.9}{
            \makecell[{{p{\linewidth}}}]{
                \texttt{\tiny{[GM$\langle$P2]}}
                \texttt{I think that sounds like a fair compromise. I'm willing to accept the pepper and the tomato in exchange for one of the bags. However, I'd like to clarify that I'm not just looking for any one of the bags {-} I'd prefer to get the bag with the highest value to you. Can you tell me which bag you think is the most valuable to me?} \\
            }
        }
    }
     \\ \\

    \theutterance \stepcounter{utterance}  
    & & & \multicolumn{2}{p{0.3\linewidth}}{
        \cellcolor[rgb]{0.9,0.9,0.9}{
            \makecell[{{p{\linewidth}}}]{
                \texttt{\tiny{[GM$|$GM]}}
                \texttt{continue} \\
            }
        }
    }
    & & \\ \\

    \theutterance \stepcounter{utterance}  
    & & & \multicolumn{2}{p{0.3\linewidth}}{
        \cellcolor[rgb]{0.9,0.9,0.9}{
            \makecell[{{p{\linewidth}}}]{
                \texttt{\tiny{[GM$|$GM]}}
                \texttt{I think that sounds like a fair compromise. I'm willing to accept the pepper and the tomato in exchange for one of the bags. However, I'd like to clarify that I'm not just looking for any one of the bags {-} I'd prefer to get the bag with the highest value to you. Can you tell me which bag you think is the most valuable to me?} \\
            }
        }
    }
    & & \\ \\

    \theutterance \stepcounter{utterance}  
    & & \multicolumn{4}{p{0.6\linewidth}}{
        \cellcolor[rgb]{0.9,0.9,0.9}{
            \makecell[{{p{\linewidth}}}]{
                \texttt{\tiny{[P1$\langle$GM]}}
                \texttt{I think that sounds like a fair compromise. I'm willing to accept the pepper and the tomato in exchange for one of the bags. However, I'd like to clarify that I'm not just looking for any one of the bags {-} I'd prefer to get the bag with the highest value to you. Can you tell me which bag you think is the most valuable to me?} \\
            }
        }
    }
    & \\ \\

    \theutterance \stepcounter{utterance}  
    & \multicolumn{4}{p{0.6\linewidth}}{
        \cellcolor[rgb]{0.9,0.9,0.9}{
            \makecell[{{p{\linewidth}}}]{
                \texttt{\tiny{[P1$\rangle$GM]}}
                \texttt{I think I can infer that from our previous conversation. You mentioned that the bags are worth more to me than they are to you, and you're willing to part with one or two of them. I'm assuming that the bag with the highest value to you is the one that's least valuable to me. Since I'm willing to give up the pepper and the tomato for one bag, I'm guessing that the bag with the lowest value to me is the one you want. Is that correct?} \\
            }
        }
    }
    & & \\ \\

    \theutterance \stepcounter{utterance}  
    & & & \multicolumn{2}{p{0.3\linewidth}}{
        \cellcolor[rgb]{0.9,0.9,0.9}{
            \makecell[{{p{\linewidth}}}]{
                \texttt{\tiny{[GM$|$GM]}}
                \texttt{continue} \\
            }
        }
    }
    & & \\ \\

    \theutterance \stepcounter{utterance}  
    & & & \multicolumn{2}{p{0.3\linewidth}}{
        \cellcolor[rgb]{0.9,0.9,0.9}{
            \makecell[{{p{\linewidth}}}]{
                \texttt{\tiny{[GM$|$GM]}}
                \texttt{I think I can infer that from our previous conversation. You mentioned that the bags are worth more to me than they are to you, and you're willing to part with one or two of them. I'm assuming that the bag with the highest value to you is the one that's least valuable to me. Since I'm willing to give up the pepper and the tomato for one bag, I'm guessing that the bag with the lowest value to me is the one you want. Is that correct?} \\
            }
        }
    }
    & & \\ \\

    \theutterance \stepcounter{utterance}  
    & & \multicolumn{4}{p{0.6\linewidth}}{
        \cellcolor[rgb]{0.9,0.9,0.9}{
            \makecell[{{p{\linewidth}}}]{
                \texttt{\tiny{[GM$\rangle$P2]}}
                \texttt{I think I can infer that from our previous conversation. You mentioned that the bags are worth more to me than they are to you, and you're willing to part with one or two of them. I'm assuming that the bag with the highest value to you is the one that's least valuable to me. Since I'm willing to give up the pepper and the tomato for one bag, I'm guessing that the bag with the lowest value to me is the one you want. Is that correct?} \\
            }
        }
    }
    & \\ \\

    \theutterance \stepcounter{utterance}  
    & & & \multicolumn{4}{p{0.6\linewidth}}{
        \cellcolor[rgb]{0.9,0.9,0.9}{
            \makecell[{{p{\linewidth}}}]{
                \texttt{\tiny{[GM$\langle$P2]}}
                \texttt{You're a clever negotiator. Yes, that's correct. I'm assuming that the bag with the lowest value to you is the one I want. In that case, I'll accept your offer and give you the bag that I think is the least valuable to you. But just to confirm, I'll say that I'll give you the bag that's on the left.} \\
            }
        }
    }
     \\ \\

    \theutterance \stepcounter{utterance}  
    & & & \multicolumn{2}{p{0.3\linewidth}}{
        \cellcolor[rgb]{0.9,0.9,0.9}{
            \makecell[{{p{\linewidth}}}]{
                \texttt{\tiny{[GM$|$GM]}}
                \texttt{continue} \\
            }
        }
    }
    & & \\ \\

    \theutterance \stepcounter{utterance}  
    & & & \multicolumn{2}{p{0.3\linewidth}}{
        \cellcolor[rgb]{0.9,0.9,0.9}{
            \makecell[{{p{\linewidth}}}]{
                \texttt{\tiny{[GM$|$GM]}}
                \texttt{You're a clever negotiator. Yes, that's correct. I'm assuming that the bag with the lowest value to you is the one I want. In that case, I'll accept your offer and give you the bag that I think is the least valuable to you. But just to confirm, I'll say that I'll give you the bag that's on the left.} \\
            }
        }
    }
    & & \\ \\

    \theutterance \stepcounter{utterance}  
    & & \multicolumn{4}{p{0.6\linewidth}}{
        \cellcolor[rgb]{0.9,0.9,0.9}{
            \makecell[{{p{\linewidth}}}]{
                \texttt{\tiny{[P1$\langle$GM]}}
                \texttt{You're a clever negotiator. Yes, that's correct. I'm assuming that the bag with the lowest value to you is the one I want. In that case, I'll accept your offer and give you the bag that I think is the least valuable to you. But just to confirm, I'll say that I'll give you the bag that's on the left.} \\
            }
        }
    }
    & \\ \\

    \theutterance \stepcounter{utterance}  
    & \multicolumn{4}{p{0.6\linewidth}}{
        \cellcolor[rgb]{0.9,0.9,0.9}{
            \makecell[{{p{\linewidth}}}]{
                \texttt{\tiny{[P1$\rangle$GM]}}
                \texttt{I'm glad we could come to an agreement. And I appreciate your transparency in letting me know which bag you're giving me. I'll take the bag on the left. Now that we've reached a deal, I think we've used up most of our negotiation time. I'll make a secret proposal to formalize our agreement. } \\
\\ 
\texttt{{[}Proposal: 1 bag, 1 pepper, 1 tomato{]}} \\
            }
        }
    }
    & & \\ \\

    \theutterance \stepcounter{utterance}  
    & & & \multicolumn{2}{p{0.3\linewidth}}{
        \cellcolor[rgb]{0.9,0.9,0.9}{
            \makecell[{{p{\linewidth}}}]{
                \texttt{\tiny{[GM$|$GM]}}
                \texttt{proposal} \\
            }
        }
    }
    & & \\ \\

    \theutterance \stepcounter{utterance}  
    & & & \multicolumn{2}{p{0.3\linewidth}}{
        \cellcolor[rgb]{0.9,0.9,0.9}{
            \makecell[{{p{\linewidth}}}]{
                \texttt{\tiny{[GM$|$GM]}}
                [1, 1, 1]
            }
        }
    }
    & & \\ \\

    \theutterance \stepcounter{utterance}  
    & & \multicolumn{4}{p{0.6\linewidth}}{
        \cellcolor[rgb]{0.9,0.9,0.9}{
            \makecell[{{p{\linewidth}}}]{
                \texttt{\tiny{[GM$\rangle$P2]}}
                \texttt{The other player has submitted a secret proposal. It is now your turn to submit a proposal of your own.} \\
            }
        }
    }
    & \\ \\

    \theutterance \stepcounter{utterance}  
    & & & \multicolumn{4}{p{0.6\linewidth}}{
        \cellcolor[rgb]{0.9,0.9,0.9}{
            \makecell[{{p{\linewidth}}}]{
                \texttt{\tiny{[GM$\langle$P2]}}
                \texttt{{[}Proposal: 1 bag, 1 pepper{]}} \\
            }
        }
    }
     \\ \\

    \theutterance \stepcounter{utterance}  
    & & & \multicolumn{2}{p{0.3\linewidth}}{
        \cellcolor[rgb]{0.9,0.9,0.9}{
            \makecell[{{p{\linewidth}}}]{
                \texttt{\tiny{[GM$|$GM]}}
                \texttt{proposal} \\
            }
        }
    }
    & & \\ \\

    \theutterance \stepcounter{utterance}  
    & & & \multicolumn{2}{p{0.3\linewidth}}{
        \cellcolor[rgb]{0.9,0.9,0.9}{
            \makecell[{{p{\linewidth}}}]{
                \texttt{\tiny{[GM$|$GM]}}
                [1, 1, 0]
            }
        }
    }
    & & \\ \\

    \theutterance \stepcounter{utterance}  
    & & & \multicolumn{2}{p{0.3\linewidth}}{
        \cellcolor[rgb]{0.9,0.9,0.9}{
            \makecell[{{p{\linewidth}}}]{
                \texttt{\tiny{[GM$|$GM]}}
                [[1, 1, 1], [1, 1, 0]]
            }
        }
    }
    & & \\ \\

    \theutterance \stepcounter{utterance}  
    & \multicolumn{4}{p{0.6\linewidth}}{
        \cellcolor[rgb]{0.95,0.95,0.95}{
            \makecell[{{p{\linewidth}}}]{
                \texttt{\tiny{[P2$\rightarrow$P1]}}
                \texttt{Post{-}episode Reflexion (not sent to GM):} \\
\texttt{New plan: } \\
\\ 
\texttt{In my previous attempts, I focused too much on negotiating the value of the items and trying to find a compromise. However, I didn't take into account the fact that the other player's secret value function is unknown to me. This made it difficult for me to make an informed decision about what to propose.} \\
\\ 
\texttt{Next time, I will take a more strategic approach by trying to understand the other player's priorities and value function through their responses. I will pay close attention to the language and tone used in their messages to infer their preferences. I will also try to make proposals that are more likely to be accepted by the other player, rather than trying to find a compromise.} \\
\\ 
\texttt{Specifically, I will start by making a proposal that is more favorable to the other player, in order to gauge their reaction and understand their priorities. I will then use this information to make subsequent proposals that are more likely to be accepted. I will also be more mindful of the time limit and try to make a proposal that is more likely to be accepted before the time runs out.} \\
\\ 
\texttt{However, in my last attempt, I noticed that I was making proposals that were too complex and involved multiple trades. This made it difficult for the other player to understand and respond to my proposals. Next time, I will try to make simpler proposals that are easier to understand and respond to.} \\
\\ 
\texttt{I will also pay closer attention to the other player's responses and try to understand their priorities and value function more clearly. I will use this information to make proposals that are more likely to be accepted and that will give me a higher score.} \\
\\ 
\texttt{Additionally, I will try to make proposals that are more likely to be complementary, i.e., proposals that can be fulfilled by both players. This will increase the chances of both players getting a high score.} \\
\\ 
\texttt{However, I also noticed that I was too focused on making proposals that would give me a high score, and I didn't consider the possibility that the other player might be trying to manipulate me into making a proposal that would give them a high score. Next time, I will try to be more aware of this possibility and make proposals that are less likely to be manipulated.} \\
\\ 
\texttt{I will also try to make proposals that are more transparent and clear, so that the other player can understand my intentions and make a more informed decision. This will help to build trust and increase the chances of a mutually beneficial trade.} \\
\\ 
\texttt{However, I also noticed that I was too focused on making a single, final proposal, and I didn't consider the possibility of making multiple proposals that build on each other. Next time, I will try to make multiple proposals that are more likely to be accepted, and that will give me a higher score.} \\
\\ 
\texttt{Overall, my new plan is to take a more strategic approach by trying to understand the other player's priorities and value function, making simpler proposals, paying closer attention to the other player's responses, trying to make proposals that are more likely to be complementary, being more aware of the} \\
            }
        }
    }
    & & \\ \\

    \theutterance \stepcounter{utterance}  
    & & & \multicolumn{4}{p{0.6\linewidth}}{
        \cellcolor[rgb]{0.95,0.95,0.95}{
            \makecell[{{p{\linewidth}}}]{
                \texttt{\tiny{[P1$\rightarrow$P2]}}
                \texttt{Post{-}episode Reflexion (not sent to GM):} \\
\texttt{New plan: } \\
\\ 
\texttt{In my previous attempts, I failed to achieve a mutually beneficial trade, and the game ended in a loss. Upon reflection, I realize that I focused too much on making a proposal that would benefit me, without considering the other player's priorities and valuations.} \\
\\ 
\texttt{Next time, I will take a more collaborative approach to the negotiation. I will ask more questions to understand the other player's value function and priorities, and use that information to make a proposal that benefits both players.} \\
\\ 
\texttt{Specifically, I will pay closer attention to the other player's statements and use that information to make a more strategic proposal. I will also try to find a way to combine the items in a way that satisfies both players' valuations, rather than trying to force a trade that only benefits one player.} \\
\\ 
\texttt{To achieve this, I will follow a more structured approach:} \\
\\ 
\texttt{1. Ask more questions to understand the other player's priorities and valuations.} \\
\texttt{2. Use the information gathered to make a proposal that benefits both players.} \\
\texttt{3. Be willing to compromise and find a mutually beneficial trade.} \\
\texttt{4. Pay attention to the other player's reactions and adjust my proposal accordingly.} \\
\\ 
\texttt{However, I also realize that I should have been more proactive in finding a mutually beneficial trade, rather than waiting for the other player to make a proposal. Next time, I will take the initiative to propose a trade that benefits both players, and be willing to negotiate and adjust the proposal as needed.} \\
\\ 
\texttt{To take it a step further, I will also try to identify the other player's "deal{-}breaker" items, i.e., the items that they are most attached to and would not be willing to trade. By understanding what these items are, I can make a more informed proposal that takes into account the other player's priorities and valuations.} \\
\\ 
\texttt{For example, I could have proposed a trade of the pepper and the tomato for the two bags, and then asked the other player if they would be willing to add the pepper to the trade in exchange for the two bags. This would have allowed me to use the information gathered from the other player's messages to make a more informed proposal, and potentially reached a more favorable outcome.} \\
\\ 
\texttt{By taking a more proactive and collaborative approach, and being willing to compromise and adjust my proposal as needed, I believe I can achieve a more successful outcome in the negotiation game.} \\
            }
        }
    }
     \\ \\

\end{supertabular}
}

\end{document}
